% Insert these subsections into Preliminary Results, replacing outdated claims.
% Uses packages already present in the supplied IEEE document.
% All values refer to trajlab_dataset_20261015_1552.zip (2026-10-15).
\subsection{Dataset Audit}
The October 15 export contains 39 scenes and 155 response records from 19
contributing participant codes: 145 drawn paths and 10 wait choices, with
three to six responses per scene. These are single drawn-path or wait
choices rather than complete rankings of a fixed candidate set. The export
contains no VLM outputs or VLM questionnaire ratings. We therefore treat
these data as an initial annotation pilot; the team must still verify which
records are study responses rather than interface tests.

Under the export's current grouping settings, 10 scenes have two geometric
route modes, nine include a wait response, and 17 contain more than one
substantial route-or-wait choice. Personal-space and group/activity tags
occur in 75 and 39 responses, respectively. These counts indicate diversity
in the recorded choices, but do not establish that every scene is
over-constrained or that every alternative route is physically feasible.

\subsection{Reference-Conditioned Drawing Pilot}
We additionally conducted an exploratory drawing exercise using GPT-6
Astra Pro in an interactive session. The assistant inspected the scene
images and existing human overlays, then authored new image-coordinate
waypoints. Shape-preserving interpolation converted the waypoints to
cubic B\'ezier curves. The exercise produced 31 movement suggestions and
eight explicit wait annotations. Because the human drawings were visible,
this is a reference-conditioned pilot, not a held-out evaluation of our
proposed in-context ranking method.

For geometric comparison, coordinates were scaled as
\begin{equation}
\phi(x,y)=\frac{(Wx,Hy)}{\sqrt{W^2+H^2}},
\end{equation}
and discrete Fr\'echet distances were computed using 64-point curves.
Distances were first averaged over human movement paths within each scene,
then averaged equally across the 31 scenes with an AI movement suggestion.
Wait outputs were not assigned geometric distances. The primary analysis
retained all non-administratively-excluded paths, including flagged
outliers and endpoint discrepancies.

\begin{table}[t]
\caption{Exploratory geometric comparison on the same 31 movement scenes.
Distances are percentages of the image diagonal; lower values indicate
closer image-space geometry, not greater safety or acceptability.}
\label{tab:reference_conditioned_pilot}
\centering
\begin{tabular}{@{}p{0.75\columnwidth}r@{}}
\toprule
\textbf{Comparison} & \textbf{Distance} \\
\midrule
AI suggestion to all human paths & 4.12\% \\
Straight start-to-goal line to human paths & 4.66\% \\
Human-to-human path pairs & 5.73\% \\
AI suggestion to nearest human path & 1.51\% \\
\bottomrule
\end{tabular}
\end{table}

The AI suggestions were closer to human drawings than the straight-line
reference in 16 of 31 scenes, with a scene-averaged distance of 4.12\%
versus 4.66\% of the image diagonal. The straight line is only a geometric
sanity check, not a collision-aware shortest-path baseline. Path-versus-wait
agreement was 120/155 responses (77.4\%); always choosing movement would
score 145/155 (93.5\%), illustrating the imbalance of the recorded choices.
The extra wait decisions have not been independently judged preferable.

\subsection{Data Quality and Limitations}
Three human drawings end more than 1\% of the image diagonal from their
scene's current goal. One \texttt{fri3} response differs by about 232 pixels
and accounts for that scene's one-member alternative route. Omitting the
three endpoint-mismatched drawings changes mean AI--human distance from
4.12\% to 4.00\%. The cause of the discrepancies must be resolved before
interpreting all geometric modes as social disagreement.

The exercise measures resemblance to visible references, not generalization,
ranking accuracy, calibrated confidence, or physical safety. No depth,
clearance, dynamic collision checking, or robot execution was performed.
The new suggestions should therefore be treated as image-space candidates
requiring independent feasibility and acceptability review. A subsequent
experiment will require frozen scene geometry, explicit waiting semantics,
recording-level train/test separation, and blinded human evaluation.
